\section{Discussion}
\label{sec:discussion}

This year's participation represents the most significant architectural evolution of our BioASQ pipeline since we began competing. Below we discuss the key findings, grounding each conclusion in the official results presented in Section~\ref{sec:results} and the internal validation experiments.

\subsection{What Worked: Retrieval}

The migration from PyTerrier PISA to pg\_textsearch for BM25 and from flat-file storage to Qdrant for dense retrieval proved successful not primarily in raw speed, but in operational reliability. Across all four batches, BM25 served as the backbone of every submission, and no index-corruption or synchronization issues were encountered---a marked improvement over previous editions where index maintenance consumed significant engineering time. The PostgreSQL-centric architecture, where documents, BM25 indexes, and metadata coexist in a single database, reduced the risk of configuration errors and simplified the submission pipeline.

Our internal validation results (Table~\ref{tab:internal-retrieval}) demonstrated that hybrid retrieval substantially outperforms BM25 alone: RRF fusion improved MRR from 0.549 to 0.696 (+27\%) and R@100 from 0.373 to 0.406 (+9\%). This pattern was mirrored in the official results, where systems using hybrid retrieval (\texttt{bioinfo-0} through \texttt{bioinfo-4}) consistently placed in the top 10--15 across multiple batches, with \texttt{bioinfo-3} achieving MAP rank 6 in Batch~1 and \texttt{bioinfo-0} reaching MAP rank 7 in Batch~3. RRF consistently outperformed weighted sum on MRR and Recall in both internal and official evaluations, confirming its robustness to the heterogeneous score distributions produced by BM25 and dense retrieval.

The reranker ensemble proved essential. Our best reranker configuration (BGE-reranker-v2-m3 applied to weighted-sum fusion) achieved an MRR of 0.738 in internal validation, substantially above the un-reranked baselines. The LLaMA Nemotron reranker achieved near-perfect validation scores (MAP-bioasq@10 of 0.9995), though we note that validation scores on 13B data do not guarantee equivalent gains on new test sets. The official ranking results validate this caution: while our systems consistently placed in the top tier, the margin between configurations was narrower than internal validation suggested. We incorporated dense retrieval results into the negative sampling process for reranker training, which we hypothesize helps the reranker generalize across both BM25 and dense retrieval paradigms; however, this specific claim awaits formal ablation testing.

Context-1 demonstrated the highest R@100 (0.453) in internal testing, suggesting strong potential as a recall-oriented first-stage method. However, its lower MRR (0.508) compared to hybrid+reranker configurations indicates that it is best used as a document source for downstream re-ranking rather than as a standalone retrieval solution.

\subsection{What Worked: Generation}

The agent quorum mechanism produced our strongest Phase~A+ and Phase~B results. In Batch~3, the quorum-based configurations achieved a perfect yes/no accuracy (rank 1) and factoid strict accuracy rank 1 (\texttt{bioinfo-0} and \texttt{bioinfo-3}). In Phase~B Batch~4, \texttt{bioinfo-3} achieved rank 4 in both factoid strict and lenient accuracy---our best factoid performance across all phases and batches.

A clear finding from our model-scale experiments was that smaller models (Gemma~4~E2B and E4B, Nemotron-3-Nano-4B) performed competitively with, and in several batches outperformed, larger models on most metrics. We attribute this to the tendency of smaller models to produce more focused representations, reducing over-elaboration and overfitting to spurious patterns in the retrieved context. The exception was list-type questions, where larger models' broader context retention proved advantageous for tracking multiple entities across documents.


A practical consideration is the computational cost of the agent quorum. Running multiple debate rounds with several large LLMs is expensive in terms of both inference cost and wall-clock time. For our largest configurations (e.g., Agents v2 with Nemotron-120B, Grok-4.1, Gemini-3-Flash, Gemini-2.5-Flash, and GPT-5-Mini, each producing multi-turn reasoning), a single question could consume substantial API credits. In many scenarios, this cost is difficult to justify relative to simpler ensembling or single-model approaches. However, the finding that lightweight local models (Gemma~4~E2B/E4B, Nemotron-3-Nano-4B) perform competitively fundamentally changes this calculus: these models can run on consumer-grade GPUs with negligible inference cost, making iterative agent loops and multi-turn debate economically viable. For tasks requiring sustained agentic reasoning—particularly where the cost of an incorrect answer is high—lightweight quorum configurations represent a practical and cost-effective solution.
The LLM-as-a-judge framework was deployed in two specific contexts: Phase~B Batch~1 (two submissions with top-$k$ answer selection at $k=3$ and $k=5$) and Phase~A+ Batch~2 (top-$k=1$ and $k=3$). The judge pipeline operates by first generating candidate answers from multiple model-prompt-context combinations (4 models $\times$ up to 6 prompts $\times$ 2 context sources), then using an LLM (Gemini~2.5~Flash) to score each answer on correctness, faithfulness, and completeness, selecting the best candidates and finally merging them with a separate model (Claude~Sonnet~4.6). While conceptually sound, the judge was used less extensively than the quorum for several reasons. The quorum's structured debate produced interpretable reasoning trails that aided debugging, and its adaptive document retention avoided the cost of generating full candidate answers from every model before evaluation. The judge remains valuable for exact answer types with objectively verifiable correctness, and a hybrid judge+quorum approach is a promising direction.

\subsection{What Needs Improvement}

List-type questions remain our weakest area across all phases. F-Measure ranks for list answers were typically in the 10--35 range in Phase~A+ and 11--41 in Phase~B, substantially behind our yes/no and factoid performance. The agent quorum's debate structure, while effective for converging on a single answer, may be suboptimal for list questions that require enumerating multiple distinct entities. Future work should explore specialized prompt strategies or agent configurations for list-type questions.

Snippet retrieval, our first attempt at this subtask, achieved modest results (best MAP rank 7 in Batch~3). The snippet sub-task is a natural bridge between retrieval and generation, and better snippets should directly improve agent quorum answer quality.

The performance gap between Phase~A+ (system-retrieved documents) and Phase~B (gold-standard documents) confirms that retrieval quality remains a bottleneck for generation. For example, factoid strict accuracy was consistently higher in Phase~B than Phase~A+ for the same system configurations, indicating that better documents directly translate to better answers. This motivates continued investment in retrieval improvements.

The variability in rankings across batches---\texttt{bioinfo-0} ranged from MAP rank 7 (B3) to 44 (B2) in Phase~A---underscores the importance of cross-batch consistency. We view robustness across batches as a more informative signal than peak single-batch performance, and future work should prioritize configurations that perform reliably across varying question distributions.

\subsection{Future Work}

Several directions emerge from this year's findings. The strong internal validation results for Context-1's recall, combined with the observation that retrieval quality directly impacts generation quality, motivate integrating Context-1 as a first-stage retrieval method in competition submissions; further experimentation is needed to evaluate its performance across more diverse question types and batch conditions. SPLADE and ColBERT remain promising retrieval approaches for future editions. For generation, dynamic prompt selection based on question type---particularly specialized handling of list questions---and the integration of the LLM judge as a post-quorum quality filter are immediate priorities. Improving snippet extraction quality should yield downstream benefits for both retrieval evaluation and agent quorum answer precision.

A key finding that warrants deeper investigation is the competitive performance of smaller language models relative to their larger counterparts in the agent quorum. Across multiple batches, the lightweight configurations (Gemma~4~E2B, E4B, and Nemotron-3-Nano-4B) performed comparably to much larger models on most metrics, with lower inference cost and latency. These results validate the use of lightweight models for the quorum and suggest that further testing with expanded evaluation protocols (covering more batches, question types, and model combinations) could refine our understanding of when and where small models suffice, potentially leading to more cost-effective submission strategies. The agent quorum itself would also benefit from more extensive testing and evaluation, particularly regarding the impact of debate round count, agent pool diversity, and convergence criteria on answer quality.
